\chapter{KẾT LUẬN}

\section{Kết quả đạt được}

\begin{itemize}
  \item Xây dựng được bộ quy tắc nghiệp vụ và danh sách thực thể bám sát chức năng thực tế của FamilyConnect.
  \item Thiết kế ERD/EER xử lý đúng điểm khó nhất của đề tài: mối kết hợp tự thân (huyết thống, hôn nhân) trên thực thể ThanhVien.
  \item Ánh xạ và chứng minh lược đồ quan hệ đạt 3NF/BCNF, loại bỏ các dị thường thêm/sửa/xóa điển hình.
  \item Đề xuất chỉ mục và phân vùng vật lý bám theo workload dự kiến của hệ thống, phù hợp SQL Server.
\end{itemize}

\section{Hạn chế và hướng phát triển}

\begin{itemize}
  \item Chưa mô hình hóa chi tiết ràng buộc thời gian không chồng lấp của HonNhan (cần trigger hoặc kiểm tra ở tầng ứng dụng).
  \item Chưa đi sâu thiết kế hạ tầng lưu trữ vector cho các dịch vụ AI (semantic search, gợi ý).
  \item Có thể mở rộng mô hình Closure Table cho cây gia phả khi số thế hệ và khối lượng truy vấn tăng cao.
  \item Ràng buộc "tối đa hai cha/mẹ trực tiếp" (2.2) và ràng buộc \{disjoint\} giữa QuanTriGiaDinh/ThanhVienThuong (3.4) chưa thể biểu diễn thuần bằng khóa chính/khóa ngoại --- cần CHECK constraint kiểu ứng dụng hoặc trigger, hiện chưa hiện thực trong phạm vi đề án.
  \item CamXuc (2.1) tạm loại khỏi phạm vi thiết kế bản này; có thể bổ sung như một thực thể yếu tương tự BinhLuan ở hướng phát triển tiếp theo.
\end{itemize}

